\subsection{Quantum Transpiling}

\subsubsection{The Quantum Computing Stack}

Before studying placement, scheduling, routing, and optimization, we need to understand \textbf{where transpiling sits inside the quantum-computing workflow}. We can summarize the whole process with five layers: (1) the quantum processor, (2) the quantum programming language, (3) the quantum algorithm, (4) the quantum compilation (transpiling), and (5) the quantum execution. The role of the \emph{transpiler} is to act as the \textbf{bridge between the abstract quantum algorithm and the physical hardware}. We will see that the \textbf{transpiler} is responsible for \textbf{taking a high-level quantum algorithm}, expressed in a quantum programming language, and \textbf{transforming it into a hardware-compatible circuit} that satisfies the constraints of a specific quantum processor.

\highspace
In general, a quantum computer is not just a chip. Many software and hardware layers cooperate to transform a high-level algorithm into physical operations executed on real qubits. Let's briefly describe the main layers of the quantum computing stack:
\begin{enumerate}
    \item \definition{Quantum Processor}. The \hl{lowest layer} is the \textbf{quantum processor}. A quantum processor is the \textbf{physical device} that \textbf{contains the qubits used for computation}. The processor is the actual \hl{hardware where quantum states evolve according to the laws of quantum mechanics.}
    \item \definition{Quantum Programming Language}. Writing microwave pulses directly would be impossible for programmers. Instead, we use a \textbf{quantum programming language} (similar to classical programming languages). A quantum programming language allows us to describe a circuit using quantum instructions such as $H$, $CNOT$, and $X$. The \hl{language hides the physical details of the hardware and provides an abstract description of the computation.}
    \item \definition{Quantum Algorithm}. A quantum algorithm is the \textbf{circuit that we want to execute}. At this stage we only think about the mathematical algorithm, without worrying about the hardware connectivity (i.e., which qubits can interact with which), gate durations (i.e., decoherence times), physical errors (i.e., noise), and qubit locations (i.e., where the qubits are physically located on the chip). The algorithm represents the \textbf{ideal computation}.
    \item \definition{Quantum Compilation} (Transpiling). This is the central topic of the chapter. A \textbf{quantum algorithm cannot usually be executed directly on hardware}. The circuit must first be adapted to the specific quantum processor. This adaptation process is called: \textbf{quantum compilation} or \definition{Transpiling}. We will see four main tasks of the transpiler: (1) \textbf{placement}, (2) \textbf{scheduling}, (3) \textbf{routing}, and (4) \textbf{optimization} (use heuristics to merge gates or rearrange operations). In general, the transpiler takes an ideal circuit and trasnforms it into another circuit that performs the same computation, respects hardware constraints, minimizes errors, and minimizes execution time.
    \item \definition{Quantum Execution}. After transpilation, the circuit is finally converted into physical control signals. It is the translation of gate-level instructions into signals sent to the processor. For superconducting qubits this means: generating microwave pulses, sending them to specific qubits, controlling interactions between qubits, performing measurements. This is the \textbf{stage where the abstract circuit becomes a real physical experiment}.
\end{enumerate}
In general, the \textbf{transpiler is the middle layer that makes the whole system work}.